% ============================================================
%  C: Como Programar -- 6ª Edição | Deitel & Deitel
%  Capítulo 10 -- Estruturas, Uniões, Manipulações de Bits
%               e Enumerações em C
%  EXERCÍCIOS (10.5 -- 10.20)
%  Abordagem incremental: Caps. 1 a 10
% ============================================================
\documentclass[12pt,a4paper]{article}
\usepackage[utf8]{inputenc}
\usepackage[T1]{fontenc}
\usepackage[brazil]{babel}
\usepackage{geometry}
\geometry{top=2.5cm,bottom=2.5cm,left=3cm,right=3cm}
\usepackage{parskip}\setlength{\parskip}{6pt}\setlength{\parindent}{0pt}
\usepackage{xcolor,tcolorbox,enumitem,listings,fancyhdr,titlesec,hyperref,booktabs,array,amsmath}
\tcbuselibrary{skins,breakable}

\definecolor{deitelblue}{RGB}{0,70,127}
\definecolor{deitelgray}{RGB}{240,242,245}
\definecolor{deitelgreen}{RGB}{0,110,60}
\definecolor{deitellightblue}{RGB}{220,235,250}
\definecolor{deitelred}{RGB}{160,20,20}
\definecolor{deitellorange}{RGB}{200,90,0}

\newtcolorbox{boaprática}[1][]{colback=deitellightblue,colframe=deitelblue,
  fonttitle=\bfseries\small,title={Boa prática de programação},breakable,#1}
\newtcolorbox{erroco}[1][]{colback=red!5,colframe=deitelred,
  fonttitle=\bfseries\small,title={Erro comum de programação},breakable,#1}
\newtcolorbox{respostabox}[1][]{colback=deitelgray,colframe=deitelblue!60,
  fonttitle=\bfseries\small\color{deitelblue},title={Resposta detalhada},breakable,#1}
\newtcolorbox{enunciadoBox}[1][]{colback=white,colframe=deitelblue,
  fonttitle=\bfseries,title={#1},breakable}
\newtcolorbox{saida}[1][]{colback=black!88,colframe=deitelblue,
  fontupper=\ttfamily\small\color{green!70},
  title={\small\color{white}Saída do programa},breakable,#1}

\setlist[enumerate,1]{label=\textbf{\alph*)},leftmargin=2em}
\setlist[itemize]{leftmargin=2em}

\lstset{
  basicstyle=\ttfamily\small,backgroundcolor=\color{deitelgray},
  frame=single,framerule=0.4pt,rulecolor=\color{deitelblue!40},
  breaklines=true,showstringspaces=false,
  keywordstyle=\color{deitelblue}\bfseries,
  commentstyle=\color{deitelgreen}\itshape,
  stringstyle=\color{deitelred},
  language=C,numbers=left,numberstyle=\tiny\color{gray},
  xleftmargin=18pt,xrightmargin=5pt,stepnumber=1,
}

\pagestyle{fancy}\fancyhf{}
\fancyhead[L]{\small\color{deitelblue}\textbf{C: Como Programar -- 6ª ed. | Deitel \& Deitel}}
\fancyhead[R]{\small\color{deitelblue}Cap. 10 -- Exercícios}
\fancyfoot[C]{\small\thepage}
\renewcommand{\headrulewidth}{0.4pt}

\titleformat{\section}[block]{\large\bfseries\color{deitelblue}}{\thesection.}{0.5em}{}[\titlerule]
\titleformat{\subsection}[block]{\normalsize\bfseries\color{deitelblue!80}}{}{}{}

\hypersetup{colorlinks=true,linkcolor=deitelblue}

\begin{document}

\begin{titlepage}
  \centering\vspace*{2cm}
  {\color{deitelblue}\rule{\linewidth}{2pt}}\\[0.4cm]
  {\huge\bfseries\color{deitelblue}C: Como Programar}\\[0.2cm]
  {\large\color{deitelblue}6ª Edição $\cdot$ Paul Deitel \& Harvey Deitel}\\[0.2cm]
  {\color{deitelblue}\rule{\linewidth}{2pt}}\\[1cm]
  {\LARGE\bfseries Capítulo 10}\\[0.3cm]
  {\Large\color{deitelblue!80}Estruturas, Uniões, Manipulações de Bits e Enumerações em C}\\[0.8cm]
  \begin{tcolorbox}[colback=deitellightblue,colframe=deitelblue,width=0.85\linewidth]
    \centering{\large\bfseries Exercícios (10.5--10.20)}\\[0.2cm]
    {\normalsize Resolvidos passo a passo, com código completo}
  \end{tcolorbox}
\end{titlepage}

\tableofcontents\newpage

% ============================================================
\section{Exercício 10.5 -- Definir Estruturas e Uniões}
% ============================================================

\begin{respostabox}
\begin{enumerate}[label=\textbf{\alph*)}]

\item \textbf{Estrutura \texttt{estoque}:}
\begin{lstlisting}
struct estoque {
    char nomePeca[ 30 ];
    int   numPeca;
    float preco;
    int   quantidade;
    int   pedido;
};
\end{lstlisting}

\item \textbf{União \texttt{dados}:}
\begin{lstlisting}
union dados {
    char   c;
    short  s;
    long   b;
    float  f;
    double d;
};
\end{lstlisting}

\item \textbf{Estrutura \texttt{endereco}:}
\begin{lstlisting}
struct endereco {
    char rua[ 25 ];
    char cidade[ 20 ];
    char estado[ 2 ];
    char cep[ 8 ];
};
\end{lstlisting}

\item \textbf{Estrutura \texttt{aluno} aninhando \texttt{endereco}:}
\begin{lstlisting}
struct aluno {
    char nome[ 15 ];
    char sobrenome[ 15 ];
    struct endereco endResid;  /* estrutura aninhada */
};
\end{lstlisting}

\item \textbf{Estrutura \texttt{teste} com 16 campos de bit:}
\begin{lstlisting}
struct teste {
    unsigned int a : 1;
    unsigned int b : 1;
    unsigned int c : 1;
    unsigned int d : 1;
    unsigned int e : 1;
    unsigned int f : 1;
    unsigned int g : 1;
    unsigned int h : 1;
    unsigned int i : 1;
    unsigned int j : 1;
    unsigned int k : 1;
    unsigned int l : 1;
    unsigned int m : 1;
    unsigned int n : 1;
    unsigned int o : 1;
    unsigned int p : 1;
};
\end{lstlisting}

\textbf{Campos de bit} permitem economizar memória quando precisamos apenas de
poucos bits para representar valores -- útil em estruturas que mapeiam registradores
de hardware ou protocolos de rede.
\end{enumerate}
\end{respostabox}

\newpage

% ============================================================
\section{Exercício 10.6 -- Acesso a Membros}
% ============================================================

\begin{respostabox}
Dadas as declarações da estrutura aninhada \texttt{cliente} e o ponteiro
\texttt{ptrCliente = \&regCliente}:

\begin{enumerate}[label=\textbf{\alph*)}]
\item \texttt{regCliente.sobrenome}
\item \texttt{ptrCliente->sobrenome}
\item \texttt{regCliente.nome}
\item \texttt{ptrCliente->nome}
\item \texttt{regCliente.numCliente}
\item \texttt{ptrCliente->numCliente}
\item \texttt{regCliente.pessoal.telefone}
\item \texttt{ptrCliente->pessoal.telefone}
\item \texttt{regCliente.pessoal.endereco}
\item \texttt{ptrCliente->pessoal.endereco}
\item \texttt{regCliente.pessoal.cidade}
\item \texttt{ptrCliente->pessoal.cidade}
\item \texttt{regCliente.pessoal.estado}
\item \texttt{ptrCliente->pessoal.estado}
\item \texttt{regCliente.pessoal.cep}
\item \texttt{ptrCliente->pessoal.cep}
\end{enumerate}

\begin{boaprática}
  \textbf{Padrão de acesso:}
  \begin{itemize}[noitemsep]
    \item Variável simples: \texttt{var.membro}
    \item Via ponteiro: \texttt{ptr->membro} (equivale a \texttt{(*ptr).membro})
    \item Aninhada: \texttt{ptr->grupo.membro} (acompanhe a hierarquia)
  \end{itemize}
\end{boaprática}
\end{respostabox}

% ============================================================
\section{Exercício 10.7 -- Embaralhamento Eficiente de Cartas}
% ============================================================

\begin{respostabox}
\begin{lstlisting}
/* Ex10.7: cartas_struct.c
   Embaralhamento e distribuicao usando struct Card */
#include <stdio.h>
#include <stdlib.h>
#include <time.h>
#define CARTAS 52

struct card {
    const char *face;
    const char *suit;
};

typedef struct card Card;

void fillDeck( Card * const baralho,
               const char * faces[],
               const char * suits[] );
void shuffleDeck( Card * const baralho );  /* Fisher-Yates */
void dealDeck( const Card * const baralho );

int main( void )
{
    /* Inicializa arrays de nomes */
    const char *faces[] = {
        "As","Dois","Tres","Quatro","Cinco","Seis","Sete",
        "Oito","Nove","Dez","Valete","Dama","Rei"
    };
    const char *suits[] = { "Copas","Ouros","Espadas","Paus" };

    Card deck[ CARTAS ];

    srand( time( NULL ) );

    fillDeck( deck, faces, suits );
    shuffleDeck( deck );
    dealDeck( deck );

    return 0;
}

void fillDeck( Card * const baralho,
               const char * faces[],
               const char * suits[] )
{
    int i;
    for ( i = 0; i < CARTAS; ++i ) {
        baralho[ i ].face = faces[ i % 13 ];
        baralho[ i ].suit = suits[ i / 13 ];
    }
}

/* Fisher-Yates: O(n), sem adiamento indefinido */
void shuffleDeck( Card * const baralho )
{
    int i, j;
    Card temp;
    for ( i = CARTAS - 1; i > 0; --i ) {
        j = rand() % ( i + 1 );
        temp = baralho[ i ];
        baralho[ i ] = baralho[ j ];
        baralho[ j ] = temp;
    }
}

/* Distribui em formato de 2 colunas */
void dealDeck( const Card * const baralho )
{
    int i;
    for ( i = 0; i < CARTAS; ++i ) {
        printf( "%-7s de %-10s%s",
                baralho[ i ].face,
                baralho[ i ].suit,
                ( i % 2 == 0 ) ? "   " : "\n" );
    }
}
\end{lstlisting}

\textbf{Pontos-chave:}
\begin{itemize}[noitemsep]
  \item \texttt{struct card} agrupa \texttt{face} e \texttt{suit} -- duas
  informações por carta.
  \item \texttt{Fisher-Yates} elimina o adiamento indefinido do algoritmo ingênuo.
  \item \texttt{Card * const baralho} -- ponteiro constante (não pode mudar para onde
  aponta) mas o conteúdo apontado pode mudar.
\end{itemize}
\end{respostabox}

\newpage

% ============================================================
\section{Exercício 10.8 -- União \texttt{integer}}
% ============================================================

\begin{respostabox}
\begin{lstlisting}
/* Ex10.8: union_integer.c
   Demonstra como diferentes tipos compartilham memoria */
#include <stdio.h>

union integer {
    char  c;
    short s;
    int   i;
    long  b;
};

int main( void )
{
    union integer u;

    /* Aceita um char */
    printf( "Digite um char: " );
    scanf( " %c", &u.c );
    printf( "\nApos atribuir char:\n" );
    printf( "  c = %c\n", u.c );
    printf( "  s = %hd\n", u.s );
    printf( "  i = %d\n", u.i );
    printf( "  b = %ld\n", u.b );

    /* Aceita um short */
    printf( "\nDigite um short: " );
    scanf( "%hd", &u.s );
    printf( "Apos atribuir short:\n" );
    printf( "  c = %c (codigo %d)\n", u.c, u.c );
    printf( "  s = %hd\n", u.s );
    printf( "  i = %d\n", u.i );
    printf( "  b = %ld\n", u.b );

    /* Aceita int e long */
    printf( "\nDigite um int: " );
    scanf( "%d", &u.i );
    printf( "Apos atribuir int:\n" );
    printf( "  c = %c (codigo %d)\n", u.c, u.c );
    printf( "  s = %hd\n", u.s );
    printf( "  i = %d\n", u.i );
    printf( "  b = %ld\n", u.b );

    printf( "\nDigite um long: " );
    scanf( "%ld", &u.b );
    printf( "Apos atribuir long:\n" );
    printf( "  c = %c (codigo %d)\n", u.c, u.c );
    printf( "  s = %hd\n", u.s );
    printf( "  i = %d\n", u.i );
    printf( "  b = %ld\n", u.b );

    return 0;
}
\end{lstlisting}

\textbf{Os valores sempre são impressos corretamente?} \textbf{NÃO!} Como todos
os membros \textit{compartilham} memória, atribuir a um membro destrói o conteúdo
dos outros. Apenas o último membro atribuído é confiável; os demais mostram
``lixo'' do ponto de vista lógico (na verdade, são reinterpretações dos mesmos
bits sob outros tipos).

\begin{boaprática}
  Uniões são úteis quando uma variável precisa ser interpretada de várias formas,
  ou quando queremos economizar memória sabendo que apenas \textit{um} dos campos
  é válido em cada momento. Em geral, mantenha um ``tag'' separado indicando qual
  membro está atualmente válido.
\end{boaprática}
\end{respostabox}

\newpage

% ============================================================
\section{Exercício 10.9 -- União \texttt{pontoFlutuante}}
% ============================================================

\begin{respostabox}
\begin{lstlisting}
/* Ex10.9: union_float.c */
#include <stdio.h>

union pontoFlutuante {
    float       f;
    double      d;
    long double x;
};

int main( void )
{
    union pontoFlutuante u;

    printf( "Digite um float: " );
    scanf( "%f", &u.f );
    printf( "  f=%f  d=%f  x=%Lf\n\n", u.f, u.d, u.x );

    printf( "Digite um double: " );
    scanf( "%lf", &u.d );
    printf( "  f=%f  d=%f  x=%Lf\n\n", u.f, u.d, u.x );

    printf( "Digite um long double: " );
    scanf( "%Lf", &u.x );
    printf( "  f=%f  d=%f  x=%Lf\n", u.f, u.d, u.x );

    return 0;
}
\end{lstlisting}

\textbf{Mesma conclusão do Ex.~10.8:} apenas o último valor atribuído faz sentido.
A diferença aqui é que os tipos \texttt{float} (4 bytes), \texttt{double} (8 bytes)
e \texttt{long double} (10--16 bytes, dependendo do sistema) têm representação
IEEE 754 muito diferentes -- interpretar os mesmos bits como tipos diferentes
produz lixo numérico.

\textbf{Tamanho da união:} a união ocupa o tamanho do \textit{maior} membro --
\texttt{sizeof(long double)} no caso.
\end{respostabox}

% ============================================================
\section{Exercício 10.10 -- Deslocamento à Direita}
% ============================================================

\begin{respostabox}
\begin{lstlisting}
/* Ex10.10: shift_direita.c */
#include <stdio.h>

void displayBits( unsigned valor );

int main( void )
{
    unsigned valor;

    printf( "Digite um inteiro nao negativo: " );
    scanf( "%u", &valor );

    printf( "\nAntes do deslocamento:  " );
    displayBits( valor );

    valor >>= 4;  /* desloca 4 bits a direita */

    printf( "Apos >>= 4:             " );
    displayBits( valor );

    return 0;
}

void displayBits( unsigned valor )
{
    unsigned i;
    unsigned mask = 1u << ( sizeof( unsigned ) * 8 - 1 );

    for ( i = 1; i <= sizeof( unsigned ) * 8; ++i ) {
        putchar( valor & mask ? '1' : '0' );
        valor <<= 1;

        if ( i % 8 == 0 ) putchar( ' ' );
    }
    putchar( '\n' );
}
\end{lstlisting}

\textbf{O que entra nos bits vagos?}
\begin{itemize}[noitemsep]
  \item Para \texttt{unsigned}: sempre 0s. Esse é chamado \textbf{deslocamento lógico}.
  \item Para \texttt{int} (signed) negativo: depende da implementação -- a maioria
  usa \textbf{deslocamento aritmético}, copiando o bit de sinal (1). Em compiladores
  modernos da família GCC, signed right shift de negativo preserva o sinal.
\end{itemize}

\begin{boaprática}
  Para resultados portáveis e previsíveis, prefira sempre tipos \texttt{unsigned}
  ao realizar manipulações de bits.
\end{boaprática}
\end{respostabox}

\newpage

% ============================================================
\section{Exercício 10.11 -- \texttt{displayBits} para 2 Bytes}
% ============================================================

\begin{respostabox}
A função \texttt{displayBits} do Ex.~10.10 já é portável: ela usa
\texttt{sizeof(unsigned) * 8} para determinar o número de bits. Em sistemas com
\texttt{unsigned} de 2 bytes (16 bits), o loop itera 16 vezes; com 4 bytes,
32 vezes.

\textbf{Versão original do livro (não portável, fixa em 4 bytes):}
\begin{lstlisting}
/* Versao fixa em 32 bits */
void displayBits( unsigned valor )
{
    int i;
    unsigned mask = 1 << 31;  /* fixo em 32 bits */

    for ( i = 1; i <= 32; ++i ) {
        putchar( valor & mask ? '1' : '0' );
        valor <<= 1;
        if ( i % 8 == 0 ) putchar( ' ' );
    }
    putchar( '\n' );
}
\end{lstlisting}

\textbf{Para 2 bytes, basta mudar 31 $\to$ 15 e 32 $\to$ 16}.
\end{respostabox}

% ============================================================
\section{Exercício 10.12 -- Multiplicação por Deslocamento}
% ============================================================

\begin{respostabox}
\textbf{Fato:} deslocar à esquerda em 1 bit equivale a multiplicar por 2.
Em \texttt{n} bits, equivale a multiplicar por $2^n$.

\begin{lstlisting}
/* Ex10.12: potencia_dois.c */
#include <stdio.h>

void displayBits( unsigned valor );

unsigned potencia2( unsigned numero, unsigned pot )
{
    return numero << pot;
}

int main( void )
{
    unsigned numero, pot;

    printf( "Digite numero e expoente: " );
    scanf( "%u%u", &numero, &pot );

    unsigned resultado = potencia2( numero, pot );

    printf( "\n%u * 2^%u = %u\n\n", numero, pot, resultado );
    printf( "Bits de %u: ", numero );    displayBits( numero );
    printf( "Bits de resultado: " );      displayBits( resultado );

    return 0;
}
\end{lstlisting}

\textbf{Por que funciona?} A representação binária ``empurra'' todos os bits para
posições mais altas. Como cada posição vale $2^k$, mover um bit uma posição à
esquerda multiplica seu peso por 2.

\textit{Exemplo:} $5 = 101_2$. $5 << 2 = 10100_2 = 20 = 5 \times 4 = 5 \times 2^2$. ✓
\end{respostabox}

\newpage

% ============================================================
\section{Exercício 10.13 -- Compactar 2 Caracteres}
% ============================================================

\begin{respostabox}
\begin{lstlisting}
/* Ex10.13: compactar.c */
#include <stdio.h>

void displayBits( unsigned valor );
unsigned compactaCaracteres( char a, char b );

int main( void )
{
    char a, b;

    printf( "Digite dois caracteres: " );
    scanf( " %c %c", &a, &b );

    printf( "\nCaractere 1 ('%c'): ", a );
    displayBits( ( unsigned )a );

    printf( "Caractere 2 ('%c'): ", b );
    displayBits( ( unsigned )b );

    unsigned compactado = compactaCaracteres( a, b );

    printf( "\nCompactado:         " );
    displayBits( compactado );

    return 0;
}

unsigned compactaCaracteres( char a, char b )
{
    unsigned resultado;

    resultado  = a;             /* coloca 'a' nos 8 bits baixos */
    resultado <<= 8;            /* desloca para os 8 bits altos */
    resultado |= ( unsigned )b; /* coloca 'b' nos 8 bits baixos */

    return resultado;
}
\end{lstlisting}

\textbf{Exemplo:} para \texttt{'A'} (65 = \texttt{01000001}) e \texttt{'B'}
(66 = \texttt{01000010}):
\begin{enumerate}[noitemsep]
  \item \texttt{resultado = 0x41 = 00000000 01000001}
  \item \texttt{resultado <<= 8} $\to$ \texttt{01000001 00000000}
  \item \texttt{resultado |= 0x42} $\to$ \texttt{01000001 01000010}
\end{enumerate}

Os dois caracteres agora ocupam 16 bits da variável \texttt{unsigned}, com \texttt{'A'}
nos bits altos e \texttt{'B'} nos bits baixos.

\textbf{Por que isso é útil?} Em sistemas embarcados e protocolos de rede, é comum
\textit{empacotar} múltiplos valores pequenos em uma variável maior para economizar
memória ou largura de banda.
\end{respostabox}

\newpage

% ============================================================
\section{Exercício 10.14 -- Descompactar 2 Caracteres}
% ============================================================

\begin{respostabox}
\begin{lstlisting}
/* Ex10.14: descompactar.c */
#include <stdio.h>

void displayBits( unsigned valor );

void descompactaCaracteres( unsigned valor, char *a, char *b )
{
    /* Mascara 65280 = 0x0000FF00 isola os bits 8-15 */
    *a = ( valor & 65280 ) >> 8;

    /* Mascara 255 = 0x000000FF isola os bits 0-7 */
    *b = valor & 255;
}

int main( void )
{
    unsigned valor;
    char a, b;

    printf( "Digite um inteiro compactado: " );
    scanf( "%u", &valor );

    printf( "\nValor compactado: " );
    displayBits( valor );

    descompactaCaracteres( valor, &a, &b );

    printf( "\nCaractere 1 ('%c'): ", a );
    displayBits( ( unsigned )a );

    printf( "Caractere 2 ('%c'): ", b );
    displayBits( ( unsigned )b );

    return 0;
}
\end{lstlisting}

\textbf{Técnica das máscaras:}
\begin{itemize}[noitemsep]
  \item \texttt{0xFF00} (65280) tem apenas os bits 8--15 acesos.
  \item \texttt{AND} com essa máscara \textit{isola} apenas esses bits.
  \item \texttt{>> 8} desloca para a posição correta de \texttt{char}.
  \item \texttt{0x00FF} (255) isola os bits 0--7 sem necessidade de deslocar.
\end{itemize}
\end{respostabox}

% ============================================================
\section{Exercício 10.15 -- Compactar 4 Caracteres}
% ============================================================

\begin{respostabox}
\begin{lstlisting}
/* Ex10.15: compactar4.c */
#include <stdio.h>

unsigned compactaCaracteres4( char a, char b, char c, char d )
{
    unsigned resultado = 0;

    resultado  = ( unsigned )a;  /* a nos bits 24-31 */
    resultado <<= 8;
    resultado |= ( unsigned )b;  /* b nos bits 16-23 */
    resultado <<= 8;
    resultado |= ( unsigned )c;  /* c nos bits 8-15 */
    resultado <<= 8;
    resultado |= ( unsigned )d;  /* d nos bits 0-7 */

    return resultado;
}

int main( void )
{
    char a, b, c, d;
    printf( "Digite 4 caracteres: " );
    scanf( " %c %c %c %c", &a, &b, &c, &d );

    unsigned compactado = compactaCaracteres4( a, b, c, d );
    printf( "Compactado: 0x%08X\n", compactado );

    return 0;
}
\end{lstlisting}

\textbf{Exemplo:} ``ABCD'' empacotado fica \texttt{0x41424344}. Este é o conceito
por trás da representação de inteiros de 4 bytes em arquivos -- e da forma como
sistemas big-endian/little-endian organizam dados na memória.
\end{respostabox}

\newpage

% ============================================================
\section{Exercício 10.16 -- Descompactar 4 Caracteres}
% ============================================================

\begin{respostabox}
\begin{lstlisting}
/* Ex10.16: descompactar4.c */
#include <stdio.h>

void descompactaCaracteres4( unsigned v,
                             char *a, char *b, char *c, char *d )
{
    unsigned mascara = 255;  /* 0x000000FF */

    *d =   v          & mascara;
    *c = ( v >>  8 )  & mascara;
    *b = ( v >> 16 )  & mascara;
    *a = ( v >> 24 )  & mascara;
}

int main( void )
{
    unsigned valor;
    char a, b, c, d;

    printf( "Digite valor compactado (em hex, ex: 0x41424344): " );
    scanf( "%x", &valor );

    descompactaCaracteres4( valor, &a, &b, &c, &d );
    printf( "Caracteres: %c %c %c %c\n", a, b, c, d );

    return 0;
}
\end{lstlisting}

\textbf{Alternativa: máscaras separadas} (técnica do enunciado):
\begin{lstlisting}
unsigned mask_d = 255u;          /* bits 0-7   */
unsigned mask_c = 255u << 8;     /* bits 8-15  */
unsigned mask_b = 255u << 16;    /* bits 16-23 */
unsigned mask_a = 255u << 24;    /* bits 24-31 */

*d =   v & mask_d;
*c = ( v & mask_c ) >> 8;
*b = ( v & mask_b ) >> 16;
*a = ( v & mask_a ) >> 24;
\end{lstlisting}

As duas abordagens são equivalentes; a primeira (deslocar antes de mascarar)
é ligeiramente mais eficiente.
\end{respostabox}

% ============================================================
\section{Exercício 10.17 -- Inverter Bits}
% ============================================================

\begin{respostabox}
\begin{lstlisting}
/* Ex10.17: reverseBits.c */
#include <stdio.h>

void displayBits( unsigned valor );

unsigned reverseBits( unsigned valor )
{
    unsigned resultado = 0;
    unsigned numBits = sizeof( unsigned ) * 8;
    unsigned i;

    for ( i = 0; i < numBits; ++i ) {
        /* Pega o bit menos significativo de valor */
        unsigned bit = valor & 1;

        /* Desloca resultado para a esquerda e adiciona o bit */
        resultado = ( resultado << 1 ) | bit;

        /* Desloca valor para o proximo bit */
        valor >>= 1;
    }

    return resultado;
}

int main( void )
{
    unsigned v;

    printf( "Digite um unsigned: " );
    scanf( "%u", &v );

    printf( "\nAntes:  " );  displayBits( v );

    unsigned inv = reverseBits( v );
    printf( "Apos:   " );    displayBits( inv );

    return 0;
}
\end{lstlisting}

\textbf{Lógica:} percorremos cada bit de \texttt{valor} (do menos para o mais
significativo) e o colocamos como bit mais significativo de \texttt{resultado}.
É como ``empilhar'' bits ao contrário.

\textit{Aplicações:} processamento de sinais digitais, FFT
(\textit{Fast Fourier Transform}), algoritmos de checksum.
\end{respostabox}

\newpage

% ============================================================
\section{Exercício 10.18 -- \texttt{displayBits} Portátil}
% ============================================================

\begin{respostabox}
\begin{lstlisting}
/* Ex10.18: displayBits_portavel.c
   Funciona em sistemas com unsigned de qualquer tamanho */
#include <stdio.h>

void displayBits( unsigned valor )
{
    /* sizeof(unsigned) * 8 = numero de bits da arquitetura */
    unsigned numBits = sizeof( unsigned ) * 8;
    unsigned mask = 1u << ( numBits - 1 );
    unsigned i;

    printf( "%-10u = ", valor );

    for ( i = 1; i <= numBits; ++i ) {
        putchar( valor & mask ? '1' : '0' );
        valor <<= 1;
        if ( i % 8 == 0 ) putchar( ' ' );
    }
    putchar( '\n' );
}

int main( void )
{
    printf( "Tamanho de unsigned: %zu bytes (%zu bits)\n",
            sizeof( unsigned ), sizeof( unsigned ) * 8 );

    displayBits( 0u );
    displayBits( 1u );
    displayBits( 255u );
    displayBits( 65535u );
    displayBits( ~0u );  /* todos os bits 1 */

    return 0;
}
\end{lstlisting}

\textbf{Chave da portabilidade:} \texttt{sizeof} é avaliado em tempo de compilação
e retorna o tamanho exato do tipo na arquitetura atual. Em qualquer sistema (2,
4, ou 8 bytes), a função funciona sem modificações.

\begin{boaprática}
  Sempre que escrever código que dependa do tamanho de tipos, use \texttt{sizeof}
  -- nunca constantes literais como 4, 32 ou 64. Seu código será automaticamente
  portátil entre arquiteturas.
\end{boaprática}
\end{respostabox}

\newpage

% ============================================================
\section{Exercício 10.19 -- Qual é o valor de X?}
% ============================================================

\begin{respostabox}
\textbf{Função suspeita:}
\begin{lstlisting}
int multiple( int num )
{
    int i, mask = 1, mult = 1;
    for ( i = 1; i <= 10; i++, mask <<= 1 ) {
        if ( ( num & mask ) != 0 ) {
            mult = 0;
            break;
        }
    }
    return mult;
}
\end{lstlisting}

\textbf{Análise passo a passo:}

A função testa os 10 bits menos significativos de \texttt{num}. Se \textit{qualquer
um} dos primeiros 10 bits for 1, retorna 0 (não é múltiplo). Caso contrário, retorna 1.

\textbf{Quando os 10 bits menos significativos são todos zero?} Quando \texttt{num}
é múltiplo de $2^{10} = 1024$.

\textbf{Verificação:}
\begin{itemize}[noitemsep]
  \item \texttt{num = 1024 = $2^{10}$}: binário \texttt{10000000000} (apenas o
  11º bit é 1). Os bits 0--9 são todos 0. Retorna 1.
  \item \texttt{num = 1025}: bit 0 é 1. Retorna 0.
  \item \texttt{num = 2048}: \texttt{100000000000} (12º bit é 1, bits 0--9 são 0).
  Retorna 1.
\end{itemize}

\textbf{Resposta: $X = 1024$ ($2^{10}$).}

A função determina se \texttt{num} é múltiplo de 1024.
\end{respostabox}

% ============================================================
\section{Exercício 10.20 -- O que \texttt{mystery} faz?}
% ============================================================

\begin{respostabox}
\textbf{Função misteriosa:}
\begin{lstlisting}
int mystery( unsigned bits )
{
    unsigned i, total = 0;
    unsigned mask = 1 << 31;

    for ( i = 1; i <= 32; i++, bits <<= 1 ) {
        if ( ( bits & mask ) == mask ) {
            total++;
        }
    }
    return !( total % 2 ) ? 1 : 0;
}
\end{lstlisting}

\textbf{Análise:}
\begin{enumerate}[noitemsep]
  \item O loop percorre os 32 bits de \texttt{bits} (assumindo 32-bit).
  \item Para cada bit, verifica se o bit \textit{mais significativo} atual é 1.
  Se for, incrementa \texttt{total}.
  \item Após o loop, \texttt{bits} foi totalmente deslocado para a esquerda.
  \item \texttt{total} agora contém o \textit{número de bits 1} (chamado
  \textit{population count} ou \textit{Hamming weight}).
  \item Retorna 1 se \texttt{total \% 2 == 0} (par), ou 0 caso contrário.
\end{enumerate}

\textbf{Conclusão:} \texttt{mystery} verifica a \textbf{paridade par} da contagem
de bits 1. Retorna 1 (verdadeiro) se há um número \textbf{par} de bits 1; retorna
0 se ímpar.

\textbf{Aplicação:} verificação de paridade é usada em transmissão de dados e
detecção de erros. Memórias ECC, RS-232, e protocolos antigos de telecomunicações
usam essa ideia para detectar erros de transmissão.
\end{respostabox}

\end{document}
